\documentclass[11pt]{article}
\usepackage{CJKutf8}
\usepackage[margin=1in]{geometry}
\usepackage{xcolor}
\usepackage{booktabs}
\usepackage{array}
\usepackage{caption}
\usepackage{graphicx}
\usepackage[hidelinks]{hyperref}

\definecolor{revblue}{RGB}{0,51,140}
\definecolor{revred}{RGB}{165,0,0}

\newcommand{\revcomment}[1]{%
  \par\vspace{0.6em}\noindent\textcolor{revblue}{\textbf{意见.}}\ %
  \textcolor{revblue}{#1}\par\vspace{0.4em}}
\newcommand{\revresponse}[1]{%
  \par\noindent\textbf{回复.}\ #1\par\vspace{0.3em}}
\newcommand{\revchange}[1]{%
  \par\noindent\textbf{拟写入正文的修改（红色）.}\par
  \noindent\textcolor{revred}{#1}\par\vspace{0.8em}}

\captionsetup{font=small,labelfont=bf,skip=4pt}
\renewcommand{\arraystretch}{1.08}

\title{\textbf{审稿意见回复}\\[0.4em]
\large SPEAR-3D: From Streaming Reconstruction to Explicit Reasoning---A Neuro-Symbolic Approach for 3D Visual Grounding}
\author{}
\date{稿件 26-3417}

\begin{document}
\begin{CJK*}{UTF8}{gbsn}
\maketitle
\thispagestyle{empty}

\noindent 尊敬的编辑与审稿人：

\medskip
\noindent 感谢对稿件的仔细阅读，以及对实验方案提出的意见。本信回复那些要求补充测量或更正已报告数字的意见。对每一条这样的意见，我们说明测的是什么、该量如何计算，以及表格应当如何阅读。准备写入正文的新表述印为红色。只涉及公式、图、文献或文字错误、且不附带新实验的意见，将在正文中修改，此处不再重复。

\bigskip
\noindent\textcolor{revblue}{\large\textbf{审稿人 1}}

\subsection*{关于重建消融}

\revcomment{重建部分看起来是在 OnlineAnySeg 上更换组件：用 SAM3 替换 CropFormer，并用标签绑定的文本嵌入替换逐裁剪 CLIP 特征。需要在 OnlineAnySeg 后端上消融 SAM3 掩码加逐裁剪 CLIP，以及所提出的标签绑定。目前看不出 AP 的提高是否来自二维骨干。}

\revresponse{我们同意，并感谢审稿人要求把两处改动分开。原稿的比较同时更换了二维分割器和语义特征，因此更高的 AP 无法单独归因于其中任何一项。我们现在在 ScanNet200 上给出 \(2\times 2\) 设计。一个因素是分割器，取 SAM3 或 CropFormer。另一个因素是语义来源，取标签绑定的 CLIP 文本或逐裁剪 CLIP 嵌入。四个方案除此之外使用同一套流式融合流程。方案 A 是本文方法。方案 C 是所引 OnlineAnySeg 的配置。差值表中固定其中一个因素再读另一列，绑定的作用和骨干的作用就可以分开看。主指标是类别无关的三维实例平均精度。AP 对 COCO 的 IoU 阈值从 0.50 到 0.95、步长 0.05 取平均。AP\(_{50}\) 与 AP\(_{25}\) 是阈值分别为 0.50 与 0.25 时的同一平均精度。分割器固定时，标签绑定在 SAM3 上把 AP 提高 2.4 个百分点，在 CropFormer 上提高 3.2 个百分点，并把 AP\(_{50}\) 分别提高 3.0 与 3.7 个百分点。语义来源固定为标签绑定文本时，骨干更换对 AP 与 AP\(_{50}\) 的影响较小，分别为 0.7 与 0.6 个百分点；两边都使用逐裁剪 CLIP 时，骨干差距更大。在 AP\(_{25}\) 上，标签绑定固定时的骨干差距为 2.6 个百分点，大于 SAM3 上的绑定差距 1.8 个百分点。因此，完整方法相对 OnlineAnySeg 的 AP\(_{25}\) 提高由两个因素共同承担，正文将不再把它写成仅由标签绑定造成。}

\revchange{类别无关的三维实例分割在 ScanNet200 上评估。AP 是 IoU 阈值从 0.50 到 0.95、步长 0.05 的 COCO 式平均精度。AP\(_{50}\) 与 AP\(_{25}\) 分别是 IoU 为 0.50 与 0.25 时的平均精度。\(2\times 2\) 表把分割器与语义来源交叉。在 AP 与 AP\(_{50}\) 上，分割器固定时，标签绑定的作用更大。在 AP\(_{25}\) 上，标签绑定固定时的骨干差距大于 SAM3 上的绑定差距，因此 AP\(_{25}\) 的提高不能只归因于标签绑定。}

\noindent 见修改稿正文表 II。

\begin{table}[ht]
\centering
\caption{ScanNet200 上的类别无关三维实例分割。每格为平均精度，单位为百分数。AP 对 IoU 阈值从 0.50 到 0.95、步长 0.05 取平均。AP\(_{50}\) 使用 IoU 0.50，AP\(_{25}\) 使用 IoU 0.25。方案 A 为本文方法。方案 C 为所引 OnlineAnySeg 的配置。方案 B 与 D 补全分割器与语义来源的交叉。}
\label{tab:recon}
\small
\begin{tabular}{llrrr}
\toprule
方案 & 配置 & AP & AP\(_{50}\) & AP\(_{25}\) \\
\midrule
A & SAM3 + 标签绑定的 CLIP 文本 & 22.5 & 40.4 & 60.5 \\
B & SAM3 + 逐裁剪 CLIP & 20.1 & 37.4 & 58.7 \\
C & CropFormer + 逐裁剪 CLIP & 18.6 & 36.1 & 53.5 \\
D & CropFormer + 标签绑定的 CLIP 文本 & 21.8 & 39.8 & 57.9 \\
\bottomrule
\end{tabular}
\end{table}

\noindent 见压缩包中的补充材料表 S10。

\begin{table}[ht]
\centering
\caption{由表~\ref{tab:recon} 直接相减得到的差值，单位是百分点。正值表示前一方案更高。“标签绑定”固定分割器，只改变语义来源。“骨干”固定语义来源，只改变分割器。}
\label{tab:recondiff}
\small
\begin{tabular}{llrrr}
\toprule
对比 & 被分离的因素 & \(\Delta\)AP & \(\Delta\)AP\(_{50}\) & \(\Delta\)AP\(_{25}\) \\
\midrule
A\(-\)B & 标签绑定，SAM3 固定 & +2.4 & +3.0 & +1.8 \\
D\(-\)C & 标签绑定，CropFormer 固定 & +3.2 & +3.7 & +4.4 \\
A\(-\)D & 骨干，标签绑定固定 & +0.7 & +0.6 & +2.6 \\
B\(-\)C & 骨干，逐裁剪固定 & +1.5 & +1.3 & +5.2 \\
A\(-\)C & 完整方法相对所引 OnlineAnySeg & +3.9 & +4.3 & +7.0 \\
\bottomrule
\end{tabular}
\end{table}

\subsection*{关于 Acc@0.3、Acc@0.5 与外部基线}

\revcomment{Acc@0.3 与 Acc@0.5 是以米为单位的质心距离阈值，而 ReferIt3D 的量是在真值候选上的选择准确率，带数值阈值的 Acc@\(k\) 则是 ScanRefer 的 IoU 约定。定位表也没有外部基线：每一行都是所提系统，只切换 NS-Grounding。至少应在这些重建上运行一种已有的零样本定位方法。此外，Acc@0.3 不可能高于 Acc@0.5，但 GLM-5 在 SR3D 上、不用 NS-Grounding 时被报告为 52.5 / 51.9。在 SR3D 上，带 NS-Grounding 的 Qwen3-plus 并不是三个骨干中最好的，正文却把 65.9 称为最好结果。}

\revresponse{感谢审稿人把原稿表中混在同一个名称下的三个量分开。质心列是所选重建实例的质心与所指物体质心之间的距离，单位为米。它不是交并比，也不是 ReferIt3D 在真值候选上的选择准确率。我们将它们改名为 Acc@0.3m 与 Acc@0.5m，并在表头标明单位。一条查询只有在两个质心都存在、距离有限时，才进入 MeanDist、Acc@0.3m 与 Acc@0.5m。Acc@0.3m 是这些查询中距离不超过 0.3\,m 的百分比，Acc@0.5m 是不超过 0.5\,m 的百分比。0.3\,m 是二者中更严的检验，因此 Acc@0.3m 不可能高于 Acc@0.5m。原稿中 GLM-5 在 SR3D 上、无 NS-Grounding 的那一格把两个数字写反了。更正后为 51.9 / 52.5。表~\ref{tab:ablation} 中其余质心格都已经满足这一顺序。我们同时加入目标实例准确率 T-Acc@0.25 与 T-Acc@0.50。在 ScanNet 真值网格上，所选实例只有在所指物体是其唯一最高点 IoU 匹配、且该 IoU 至少为 0.25 或 0.50 时才算正确。没有选择的查询算错误。质心落在 0.5\,m 以内的邻近物体，只要另一物体的 IoU 更高，仍然算错误。这是质心准确率所不做的身份检验。约束满足方法 CSVG 与本文各行使用同一套流式重建，因此外部比较的几何是同一套。在 NR3D 上，CSVG 的 Acc@0.5m 为 47.7，带 NS-Grounding 的 Qwen3-plus 为 58.8。在 SR3D 上，CSVG 为 47.0，带 NS-Grounding 的 GLM-5 为 67.0。消融表中，NR3D 上最高的 Acc@0.5m 是带 NS-Grounding 的 Qwen3-plus（58.8）。SR3D 上最高的 Acc@0.5m 是带 NS-Grounding 的 GLM-5（67.0）。带 NS-Grounding 的 Qwen3-plus 在 SR3D 上为 65.9，因此正文将不再把 65.9 称为 SR3D 的最好结果。两个数据集使用同一套重建。SR3D 的语句由空间关系模板生成，NR3D 的语句是自由的人工描述，所以 NR3D 较低反映的是语言，而不是更弱的重建。表~\ref{tab:ablation} 的每一行中，打开 NS-Grounding 都会提高 Acc@0.5m 和两个目标实例准确率。}

\revchange{MeanDist 是所选重建实例质心与所指物体质心之间的平均欧氏距离，单位为米，只对两个质心都存在的查询取平均。Acc@0.3m 与 Acc@0.5m 是这同一批查询中，距离不超过 0.3\,m 与 0.5\,m 的百分比。没有选出实例的查询不进入这三列。目标实例准确率不采用这一省略：没有选择即算错误。T-Acc@0.25 与 T-Acc@0.50 是已评分查询中的百分比：在真值网格上，所指的 ScanNet 物体必须是所选实例的唯一最高点 IoU 匹配，且该 IoU 至少为 0.25 或 0.50。质心更近的邻近物体，只要另一物体的 IoU 更高，就算错误。最高的 Acc@0.5m 按数据集分别报告。}

\noindent 见修改稿正文表 III。

\begin{table}[ht]
\centering
\caption{在本文流式重建上的定位。CSVG 是约束满足基线，与本文各行使用同一套重建。带 NS-Grounding 的行是表~\ref{tab:ablation} 中该数据集 Acc@0.5m 最高的骨干：NR3D 为 Qwen3-plus，SR3D 为 GLM-5。MeanDist 单位为米，越低越好。Acc@0.3m 与 Acc@0.5m 是质心距离准确率，单位为百分数，分母是具有有限质心距离的查询。T-Acc@0.25 与 T-Acc@0.50 是目标实例准确率，单位为百分数，分母是已评分查询，没有选择算错误。}
\label{tab:csvg}
\small
\begin{tabular}{llrrrrr}
\toprule
方法 & 数据 & MeanDist & Acc@0.3m & Acc@0.5m & T-Acc@0.25 & T-Acc@0.50 \\
\midrule
CSVG & NR3D & 1.10 & 45.2 & 47.7 & 41.3 & 33.9 \\
Qwen3-plus + NS & NR3D & 1.02 & 56.3 & 58.8 & 50.4 & 40.8 \\
CSVG & SR3D & 1.34 & 46.1 & 47.0 & 43.8 & 37.1 \\
GLM-5 + NS & SR3D & 0.97 & 65.6 & 67.0 & 63.2 & 53.8 \\
\bottomrule
\end{tabular}
\end{table}

\noindent 见压缩包中的补充材料表 S1。

\begin{table}[ht]
\centering
\caption{更正后的定位消融。NS 表示 NS-Grounding；无 NS 时由语言模型直接选择实例。MeanDist、Acc@0.3m 与 Acc@0.5m 的定义同表~\ref{tab:csvg}。T-Acc@0.25 与 T-Acc@0.50 是两个 IoU 阈值下的目标实例准确率。GLM-5 在 SR3D 上、无 NS-Grounding 的一行为 51.9 / 52.5，因此 Acc@0.3m 不高于 Acc@0.5m。各数据集中最高的 Acc@0.5m 分别是 NR3D 上带 NS 的 Qwen3-plus（58.8）和 SR3D 上带 NS 的 GLM-5（67.0）。}
\label{tab:ablation}
\resizebox{\linewidth}{!}{%
\begin{tabular}{llrrrrrrrrrr}
\toprule
 & & \multicolumn{5}{c}{NR3D} & \multicolumn{5}{c}{SR3D} \\
\cmidrule(lr){3-7}\cmidrule(lr){8-12}
语言模型 & NS & 距离 & @0.3m & @0.5m & T@.25 & T@.50 & 距离 & @0.3m & @0.5m & T@.25 & T@.50 \\
\midrule
Qwen3-plus & 无 & 1.36 & 38.7 & 43.5 & 33.2 & 28.2 & 1.32 & 53.8 & 54.8 & 37.2 & 31.1 \\
Qwen3-plus & 有 & 1.02 & 56.3 & \textbf{58.8} & 50.4 & 40.8 & 0.96 & 64.7 & 65.9 & 61.4 & 52.6 \\
deepseek-v3.2 & 无 & 1.33 & 42.0 & 44.6 & 30.9 & 26.0 & 1.24 & 53.1 & 54.1 & 36.3 & 29.7 \\
deepseek-v3.2 & 有 & 1.27 & 51.1 & 54.3 & 46.0 & 39.4 & 1.00 & 64.8 & 66.0 & 64.0 & 57.0 \\
GLM-5 & 无 & 1.35 & 37.8 & 41.5 & 32.3 & 29.1 & 1.28 & 51.9 & 52.5 & 35.7 & 29.5 \\
GLM-5 & 有 & 1.18 & 50.2 & 53.4 & 44.4 & 36.2 & 0.97 & 65.6 & \textbf{67.0} & 63.2 & 53.8 \\
\bottomrule
\end{tabular}}
\end{table}

\subsection*{关于速度}

\revcomment{OnlineAnySeg 报告 15 FPS，而稿件表 I 约为 4。请给出硬件、关键帧间隔，以及分阶段耗时。SAM3 比 CropFormer 更重。10 FPS 的说法并不独立于提示词数量。}

\revresponse{我们同意，原稿把不止一种速率都称为 FPS。下面的计时在 NVIDIA RTX 3090 上测量，即正文所写的 GPU。OnlineAnySeg 发表的约 15 FPS 是在 RTX 4090 上测量的，并且按我们对该报告的理解，它对应的阶段窄于完整的 RGB-D 流。因此它与表~\ref{tab:fps} 的两列都不是同一个量。我们的端到端速率统计的是一个 ScanNet 场景上每秒处理的 RGB-D 帧数。SAM3 加标签绑定特征的平均为 10.92 帧/秒，CropFormer 加逐裁剪 CLIP 特征的平均为 5.32 帧/秒。这就是正文中约 10 FPS 所对应的速率。它不是分割调用速率。两个方案都是每 10 个关键帧调用一次分割。一次 SAM3 调用耗时 241\,ms（4.16 次/秒）。一次 CropFormer 调用耗时 1099\,ms（0.91 次/秒），约为前者的 4.6 倍。若只比较单次调用的墙钟时间，SAM3 更重；端到端流却是 SAM3 方案更快，因为标签绑定特征不必再做逐裁剪嵌入，而且分割并不在每一帧上运行。融合在 SAM3 方案上为 330\,ms，在 CropFormer 方案上为 453\,ms。合并分别为 487\,ms 与 1307\,ms。时间旁边的速率是该时间的倒数，保留两位小数，只表示该阶段每秒的调用次数、帧数或合并次数。它们都不是端到端 FPS。CropFormer 的合并速率 0.77 次/秒也不是发表的 15 FPS。词表扫描只改变一次分割帧所看到的提示词列表。延迟是该帧的墙钟时间。帧率是 \(1000\) 除以该延迟，因此仍然是分割帧速率。\(V{=}10\) 时为 3.22 帧/秒，\(V{=}100\) 时为 0.61 帧/秒。这一扫描中的设备内存峰值保持在约 16\,GB（从 16105\,MB 到 16084\,MB），而每帧平均片段数从 8.3 升到 16.8。因此，单独一句 10 FPS 描述的是所提方案的端到端流，并不能描述每一个词表规模。}

\revchange{端到端 FPS 是单个场景上每秒处理的 RGB-D 帧数，在 RTX 3090 上测量。分割每 10 个关键帧调用一次。分割时间是一次分割调用的墙钟时间，分割吞吐是该时间的倒数。融合与合并时间是这两个阶段的墙钟时间，旁边的速率是这些时间的倒数。词表扫描报告的是分割帧延迟，不是端到端 FPS。约 10 FPS 的端到端数字指的是所提方案，并不声称对每一个词表规模成立。}

\noindent 见压缩包中的补充材料表 S3。

\begin{table}[ht]
\centering
\caption{ScanNet 场景上的端到端流式速率，单位为帧/秒，在 RTX 3090 上测量。方案 A 为 SAM3 加标签绑定特征。方案 C 为 CropFormer 加逐裁剪 CLIP 特征。最后一行是上列场景的算术平均。该速率统计的是已处理的 RGB-D 帧，不是分割调用速率。}
\label{tab:fps}
\small
\begin{tabular}{lrr}
\toprule
场景 & 方案 A（帧/秒） & 方案 C（帧/秒） \\
\midrule
scene0432\_00 & 11.42 & 5.50 \\
scene0527\_00 & 10.71 & 4.82 \\
scene0559\_00 & 10.62 & 4.41 \\
scene0494\_00 & 10.82 & 6.60 \\
scene0689\_00 & 11.04 & 5.25 \\
平均 & 10.92 & 5.32 \\
\bottomrule
\end{tabular}
\end{table}

\noindent 见压缩包中的补充材料表 S4。

\begin{table}[ht]
\centering
\caption{RTX 3090 上的分阶段墙钟时间。每格中的速率是所给墙钟时间的倒数，保留两位小数。分割每 10 个关键帧调用一次，因此其调用速率不是表~\ref{tab:fps} 的端到端帧率。融合与合并的速率同样只描述该阶段。}
\label{tab:stages}
\small
\begin{tabular}{lll}
\toprule
阶段 & SAM3 + 标签绑定 & CropFormer + 逐裁剪 \\
\midrule
分割 & 241\,ms/次（4.16 次/秒） & 1099\,ms/次（0.91 次/秒） \\
融合 & 330\,ms（3.03 帧/秒） & 453\,ms（2.21 帧/秒） \\
合并 & 487\,ms（2.05 次/秒） & 1307\,ms（0.77 次/秒） \\
\bottomrule
\end{tabular}
\end{table}

\noindent 见压缩包中的补充材料表 S5。

\begin{table}[ht]
\centering
\caption{提示词表规模 \(V\) 增大时的分割帧开销，在 RTX 3090 上测量。延迟是一次分割帧的墙钟时间，单位为毫秒。帧率是 \(1000\) 除以该延迟，不是端到端 FPS。峰值内存是设备显存的高水位，单位为 MB。最后一列是每帧产生的平均片段数。}
\label{tab:vocab}
\small
\begin{tabular}{rrrrr}
\toprule
\(V\) & 延迟（ms） & 帧率 & 峰值内存（MB） & 平均片段数 \\
\midrule
10 & 311 & 3.22 & 16105 & 8.3 \\
30 & 509 & 1.97 & 16098 & 11.3 \\
50 & 762 & 1.31 & 16095 & 12.1 \\
100 & 1627 & 0.61 & 16084 & 16.8 \\
\bottomrule
\end{tabular}
\end{table}

\subsection*{关于 SceneNN}

\revcomment{在 SceneNN 上，本文方法不如 OnlineAnySeg，正文把原因归于标注不完整，但没有证据。若提示词表就是类别清单，缺失类别是这一设计的结果。}

\revresponse{我们同意，原稿在没有测量的情况下断言标注不完整，并感谢审稿人把这一说法连接到提示词表。提示词是一份固定的类别名清单。真值实例的类别若属于这些名称，则在词表内；其余实例在词表外。在 SceneNN 上，61.3\% 的真值实例在词表内，38.7\% 不在，两行占比之和为 100\%。单个标签 \texttt{otherprops} 占全部真值实例的 33.8\%，是词表外份额的主体。召回率是在给定 IoU 下被匹配到的真值实例比例，在词表内和词表外分别计算。词表内召回率在 IoU 0.50 时为 56.7\%，在 IoU 0.25 时为 85.1\%。词表外召回率在同样两个阈值下为 37.8\% 与 61.3\%，因此未命名的实例仍可以发生几何匹配。类别匹配的顶点覆盖是更严的量：只有预测类别与真值类别一致的顶点才计入。该覆盖在词表内为 61.1\%，在词表外为 0.0\%，因为没有任何提示词命名缺失的类别，匹配无法被计为类别正确。另外，完全没有标签的真值顶点，按场景平均占 28.2\%。相对一个能够输出我们词表之外类别的方法，这一差距是提示词清单的性质。修改后的正文将报告这些测得的比例，而不再使用未经量化的标注评论。}

\revchange{在 SceneNN 上，61.3\% 的真值实例属于提示词表中的类别，38.7\% 不属于。其中 \texttt{otherprops} 一类占全部实例的 33.8\%。词表内召回率在 IoU 0.50 / 0.25 下为 56.7\% / 85.1\%，词表外为 37.8\% / 61.3\%。类别匹配的顶点覆盖在词表内为 61.1\%，在词表外为 0.0\%。平均每个场景有 28.2\% 的真值顶点没有标签。词表外覆盖为零，是因为提示词清单中没有该类：不是提示词的类别不能被计为类别匹配。}

\noindent 见压缩包中的补充材料表 S6。

\begin{table}[ht]
\centering
\caption{SceneNN 的词表覆盖。“词表内”指真值类别是提示词之一。召回率是在给定 IoU 下被匹配上的真值实例比例。类别匹配顶点覆盖是预测类别与真值类别一致的顶点比例。两行实例占比之和为 100\%。}
\label{tab:scenenn}
\small
\begin{tabular}{lr}
\toprule
量 & 数值 \\
\midrule
属于词表的真值实例 & 61.3\% \\
不属于词表的真值实例 & 38.7\% \\
全部实例中标记为 \texttt{otherprops} 的比例 & 33.8\% \\
词表内召回，IoU 0.50 / 0.25 & 56.7\% / 85.1\% \\
词表外召回，IoU 0.50 / 0.25 & 37.8\% / 61.3\% \\
词表内的类别匹配顶点覆盖 & 61.1\% \\
词表外的类别匹配顶点覆盖 & 0.0\% \\
无标签真值顶点，按场景平均 & 28.2\% \\
\bottomrule
\end{tabular}
\end{table}

\bigskip
\noindent\textcolor{revblue}{\large\textbf{审稿人 3}}

\subsection*{关于标签绑定与 OVI-MAP}

\revcomment{把语义类别写入三维融合已经很常见。请说明标签介导的特征绑定，与地图建好之后再绑定类别有何不同，并与 OVI-MAP 的语义 TSDF 比较。}

\revresponse{感谢审稿人追问类别在流程中的位置。这一点我们原先留在字面之外。标签绑定在跨帧关联之前，就把提示词名称的嵌入写到二维掩码上。类别因此是融合的输入：它参与决定哪些观测属于同一个三维实例。建图后的分类器正好相反。几何实例已经存在，类别事后由一次单独的查询选出。为了固定掩码，我们取一套 SAM3 重建，在同一批实例上比较两种赋标签方式。标签绑定使用融合前写入的提示词文本，地图存在之后不再发出类别查询，因此最后一列为 0。事后分类对每个对齐实例，在 200 个类别名上做一次 CLIP argmax，因此最后一列为 1。mIoU 是各类上预测标签与真值标签交并比的平均。mAcc 是各类准确率的平均。二者都是 OVI-MAP 所用按类 IoU 评估返回的分数。在这一配对比较中，标签绑定的 mIoU 为 0.2925、mAcc 为 0.3766，事后赋值为 0.1532 与 0.2596。表~\ref{tab:ovimap} 把同一类分数写成百分数，放在 OVI-MAP 表 3 已发表数字的旁边。那一行是引用，并未重跑：OVI-MAP 在地图建好之后于视图上查询 SigLIP。本文一行在融合前绑定提示词名称。两行的场景、分割器和语义模型也不同。该表比较的是机制和已发表的工作点，是文献对照，不是在同一套重建上的配对胜负。}

\revchange{标签绑定在跨帧关联之前把提示词文本嵌入写到掩码上，因此类别参与融合，地图存在之后不再发出类别查询。事后分类不改变实例几何，然后在类别名列表上做一次 CLIP argmax。mIoU 与 mAcc 是 OVI-MAP 所用按类 IoU 评估给出的按类交并比平均和按类准确率平均。表~\ref{tab:bind} 在一套共享的 SAM3 重建上以分数报告它们。表~\ref{tab:ovimap} 把同一类分数写成百分数，并与所引 OVI-MAP 表 3 并列。这两行的场景、分割器和语义来源不同。}

\noindent 见压缩包中的补充材料表 S7。

\begin{table}[ht]
\centering
\caption{一套 SAM3 重建上的语义标签。两行使用相同的实例掩码，因此比较只改变类别何时被赋值。mIoU 是按类交并比的平均，mAcc 是按类准确率的平均，二者都是 \([0,1]\) 上的分数。最后一列是地图已经存在之后、每个实例的类别查询次数。}
\label{tab:bind}
\small
\begin{tabular}{lrrr}
\toprule
语义赋值 & mIoU & mAcc & 每实例查询次数 \\
\midrule
融合前绑定的提示词文本 & 0.2925 & 0.3766 & 0 \\
事后裁剪 CLIP，在 200 个类别名上取 argmax & 0.1532 & 0.2596 & 1 \\
\bottomrule
\end{tabular}
\end{table}

\noindent 见压缩包中的补充材料表 S8。

\begin{table}[ht]
\centering
\caption{按 OVI-MAP 表 3 的写法，以百分数给出的语义分割分数。OVI-MAP 一行引自该表，并未重跑；SigLIP 在建图后于视图上查询。SPEAR-3D 一行是本文的按类评估，提示词名称在融合前绑定。场景集合、分割器和语义来源不同，因此两行是文献对照，而不是同一套重建上的配对比较。}
\label{tab:ovimap}
\small
\begin{tabular}{llrr}
\toprule
方法 & 语义来源 & mIoU & mAcc \\
\midrule
OVI-MAP（引用表 3） & 建图后的 SigLIP & 17.5 & 27.6 \\
SPEAR-3D（本次评估） & 融合前绑定的提示词名称 & 30.6 & 39.8 \\
\bottomrule
\end{tabular}
\end{table}

\subsection*{关于七个关系打分器}

\revcomment{第 III-B.4 节引入七个手工打分器，既没有说明这个集合是否充分，也没有对每个打分器做消融。}

\revresponse{感谢审稿人追问每个打分器是否有贡献。这七个打分器是接触、方向、竖直、between、次序、侧面和邻近。每一个都是候选实例与语句所指物体之间的几何检验。我们并不声称这份清单穷尽了空间语言。能够给出的是：在其余地面器固定时，每一项检验的边际贡献。地面器给出排名第一的实例。消融每次去掉一个打分器，然后再把七个全部去掉。本表的指标是目标实例准确率。一次查询在所指物体是所选重建实例于 ScanNet 真值网格上的唯一最高点 IoU、且该 IoU 达到相应阈值时算正确。缺少选择算错误。T-Acc@0.25 的阈值是 0.25，T-Acc@0.50 的阈值是 0.50。表~\ref{tab:scorers} 中的百分数只与本表其他行比较。完整集合的 T-Acc@0.25 为 50.4，T-Acc@0.50 为 40.8，MeanDist 为 0.914\,m。去掉接触后下降最多，两项分别为 44.3（\(-6.1\) 个百分点）和 37.5（\(-3.3\) 个百分点），MeanDist 升至 1.050\,m。方向使两项下降 3.6 和 2.0 个百分点，到 46.8 和 38.8。侧面使两项下降 2.8 和 2.0 个百分点，到 47.6 和 38.8。去掉 between 使两项下降 2.4 和 4.4 个百分点，到 48.0 和 36.4。竖直使两项下降 1.5 和 2.0 个百分点，到 48.9 和 38.8。次序使两项下降 1.5 和 0.7 个百分点，到 48.9 和 40.1。去掉邻近使两项下降 1.3 和 0.6 个百分点，到 49.1 和 40.2，MeanDist 为 0.932\,m。七个全部去掉后为 33.2 和 28.2，即 \(50.4-33.2=17.2\) 与 \(40.8-28.2=12.6\) 个百分点，报告为 \(-17.2\) 和 \(-12.6\)，MeanDist 为 1.472\,m。每一次去掉都会降低两档 T-Acc。联合下降大于任何一次单独下降，说明这个集合是在共同起作用，而不是由某一个打分器承载其余各项。}

\revchange{T-Acc@0.25 与 T-Acc@0.50 是所指物体成为所选实例唯一最高点 IoU、且 IoU 分别达到 0.25 与 0.50 的查询所占的百分比。缺少选择算错误。每个 \(\Delta\) 都是相对完整集合的差值。MeanDist 是具有有限距离的查询上的平均质心距离，单位为米。}

\noindent 见修改稿正文表 IV。

\begin{table}[ht]
\centering
\caption{七个关系打分器的留一消融。T-Acc@0.25 与 T-Acc@0.50 的单位为百分数。所指物体是所选实例在相应 IoU 阈值上的唯一最高点匹配时，该查询算正确。缺少选择算错误。每个 \(\Delta\) 都是相对完整集合的差值。MeanDist 单位为米。“七个都去掉”一行是 \(50.4-33.2=17.2\) 与 \(40.8-28.2=12.6\)。}
\label{tab:scorers}
\small
\begin{tabular}{lrrrrr}
\toprule
配置 & T-Acc@0.25（\%） & \(\Delta_{0.25}\) & T-Acc@0.50（\%） & \(\Delta_{0.50}\) & MeanDist（m） \\
\midrule
完整集合 & 50.4 & --- & 40.8 & --- & 0.914 \\
去掉接触 & 44.3 & \(-6.1\) & 37.5 & \(-3.3\) & 1.050 \\
去掉方向 & 46.8 & \(-3.6\) & 38.8 & \(-2.0\) & 1.024 \\
去掉侧面 & 47.6 & \(-2.8\) & 38.8 & \(-2.0\) & 0.922 \\
去掉 between & 48.0 & \(-2.4\) & 36.4 & \(-4.4\) & 1.034 \\
去掉竖直 & 48.9 & \(-1.5\) & 38.8 & \(-2.0\) & 0.997 \\
去掉次序 & 48.9 & \(-1.5\) & 40.1 & \(-0.7\) & 0.968 \\
去掉邻近 & 49.1 & \(-1.3\) & 40.2 & \(-0.6\) & 0.932 \\
七个都去掉 & 33.2 & \(-17.2\) & 28.2 & \(-12.6\) & 1.472 \\
\bottomrule
\end{tabular}
\end{table}

\subsection*{关于传感器退化}

\revcomment{全部实验都是离线的。没有物理机器人演示，没有传感器噪声和运动模糊下的评估，也没有使用定位输出的闭环任务。}

\revresponse{感谢审稿人追问重建在输入流退化时的表现。本次修改补充运动模糊、深度噪声和丢帧的受控合成实验，评分使用表~\ref{tab:recon} 所定义的同一套类别无关 AP、AP\(_{50}\) 与 AP\(_{25}\)。本次修改不加入闭环控制器。本文的实验和创新点是流式重建与语言定位，不是抓取，也不是闭环抓取控制。本次修改不加入新的物理机器人定量实验。演示视频放在压缩包里的补充材料中，节名是“机器人演示视频”，该节的文件位置留空，便于放入视频。补充材料 PDF 和视频文件放在压缩包中，不放进正文 PDF。见压缩包中的这一节。合成实验是我们能够用论文中已有重建指标定义的测量。模糊和深度噪声作用在与干净基线相同的均匀子采样帧上，这些帧从每条完整序列中均匀抽取。模糊只改变 RGB。其方向是相邻帧之间相机平移在图像上的投影：平移变换到相机坐标系后，核的角度为 \(\mathrm{atan2}(\Delta y/\Delta z,\Delta x/\Delta z)\)。核是长度为 7、21 或 41 像素的均匀线段，长度为奇数，用线性滤波施加。位姿无效时退化为水平核。深度噪声只改变深度。对以米为单位的有效深度 \(d\)，噪声标准差为 \(\sigma_{\mathrm{mm}}=0.5+b\,d^{1.5}\)，其中 \(b\in\{0.005,0.02,0.04\}\)。在 4\,m 处，这些值约为 0.54、0.66 和 0.82\,mm，因此高斯项本身低于一毫米。噪声各行上更大的变化是空洞率：分别把 3\%、10\% 或 20\% 的有效深度像素置零，并与三个 \(b\) 配对。丢帧不改被保留的帧。相对干净基线，从完整序列中均匀保留 133、100 或 67 帧。干净 AP 为 21.4 / 39.5 / 61.7。模糊和深度噪声并不单调下降：强模糊为 22.6 / 42.0 / 63.2，高于干净行，严重模糊又回到接近干净的 AP。丢帧则单调下降，从保留 133 帧时的 19.3 / 34.8 / 54.4 降到保留 67 帧时的 13.2 / 27.6 / 42.9。因此该指标能够检出视角减少。在本实验中，它并不对每一种模糊或噪声强度都给出单调惩罚。}

\revchange{运动模糊、深度噪声和丢帧施加在 ScanNet 的 RGB-D 流上，并用重建表中的类别无关 AP 评分。模糊只改变 RGB，核为长度 7、21 或 41 像素的均匀线段，方向由相机平移决定。深度噪声使用 \(\sigma_{\mathrm{mm}}=0.5+b\,d^{1.5}\)，并配合 3\%、10\% 或 20\% 的空洞率。丢帧从完整序列中均匀保留一个子集，且不修改被保留的帧。保留的帧越少，AP 越低。AP 并不随模糊或噪声强度单调下降。}

\noindent 见压缩包中的补充材料表 S9。

\begin{table}[ht]
\centering
\caption{ScanNet 场景上的合成退化。AP、AP\(_{50}\) 与 AP\(_{25}\) 是表~\ref{tab:recon} 所定义的类别无关实例分割分数。模糊和噪声复用干净基线的均匀子采样帧。丢帧在完整序列上重新采样，且不修改被保留的帧。高斯深度噪声在 4\,m 处低于一毫米；噪声各行上更大的变化是空洞率。}
\label{tab:robust}
\small
\begin{tabular}{llrrr}
\toprule
条件 & 生成方式 & AP & AP\(_{50}\) & AP\(_{25}\) \\
\midrule
干净 & 均匀子采样；RGB 与深度不变 & 21.4 & 39.5 & 61.7 \\
弱模糊 & 运动模糊核，7\,px，只改 RGB & 21.2 & 38.5 & 60.9 \\
强模糊 & 运动模糊核，21\,px，只改 RGB & 22.6 & 42.0 & 63.2 \\
严重模糊 & 运动模糊核，41\,px，只改 RGB & 21.7 & 40.7 & 59.7 \\
弱深度噪声 & \(b{=}0.005\)，3\% 深度空洞 & 21.9 & 40.3 & 63.2 \\
强深度噪声 & \(b{=}0.02\)，10\% 深度空洞 & 22.0 & 39.9 & 63.4 \\
严重深度噪声 & \(b{=}0.04\)，20\% 深度空洞 & 21.5 & 38.6 & 61.7 \\
133 帧 & 在完整序列上均匀保留 133 帧 & 19.3 & 34.8 & 54.4 \\
100 帧 & 在完整序列上均匀保留 100 帧 & 17.0 & 32.6 & 50.0 \\
67 帧 & 在完整序列上均匀保留 67 帧 & 13.2 & 27.6 & 42.9 \\
\bottomrule
\end{tabular}
\end{table}

\bigskip
\noindent\textcolor{revblue}{\large\textbf{审稿人 5}}

\subsection*{关于不一致的数字与标准定位指标}

\revcomment{在表 II 中，GLM-5 在 SR3D 上、不使用所提定位模块时，Acc@0.3\,=\,52.5 而 Acc@0.5\,=\,51.9。若二者都是质心距离，这不可能出现。称某一个模型最好与表格不符。请复查评估和全部报告数字。也请报告 NR3D 与 SR3D 所用的标准目标实例准确率，说明查询过滤，并与结构化或基于约束的定位方法比较，而不只与直接的语言模型提示比较。附近的错误物体必须算错。}

\revresponse{我们复查了每一个质心单元格，并感谢审稿人指出这一对调。Acc@0.3m 与 Acc@0.5m 都是质心距离。唯一违反 Acc@0.3m \(\le\) Acc@0.5m 的是 GLM-5 在 SR3D 上、无 NS-Grounding 的一格，原稿印为 52.5 / 51.9。两个数字已经对调。更正后为 51.9 / 52.5。表~\ref{tab:ablation} 的其余单元格都已经满足这一顺序，且任一单元格中 T-Acc@0.50 都不高于 T-Acc@0.25。“某一个模型最好”现在按数据集、按列陈述。Acc@0.5m 在 NR3D 上的最高值是带 NS-Grounding 的 Qwen3-plus（58.8），在 SR3D 上的最高值是带 NS-Grounding 的 GLM-5（67.0）。摘要和结论将使用这两句，不再把 65.9 称为 SR3D 的最好结果。直接的语言模型提示是表~\ref{tab:ablation} 中“无 NS”的一组：模型在没有关系打分器的情况下选择实例。CSVG 是另外的基于约束的方法，在同一套重建上评估，见表~\ref{tab:csvg}。目标实例准确率是新增的两列。它在 ScanNet 真值网格上计算。所选重建实例只有在其最高点 IoU 就是所指物体、该 IoU 达到阈值、且没有其他物体与这一最高 IoU 并列时才算正确。查询过滤因列而异，这也是质心成功仍可能成为目标实例失败的原因。MeanDist、Acc@0.3m 与 Acc@0.5m 只使用距离有限的查询，因此没有产生实例的查询不进入分母。T-Acc 不采用这一省略：没有选择即算错误，分母是场景网格与所指物体都能够评分的查询。落在所指质心 0.5\,m 以内的错误物体，只要另一物体的点 IoU 更高，T-Acc 仍为零。表~\ref{tab:scorers} 对留一打分器的选择使用同一个 T-Acc 定义，其百分数只在该表内部比较。}

\revchange{定位一节将写明：Acc@0.3m 与 Acc@0.5m 是质心距离，其分母只保留距离有限的查询；T-Acc@0.25 与 T-Acc@0.50 是真值网格上的目标实例准确率，没有选择算错误。对 T-Acc 而言，只要另一物体的点 IoU 更高，附近的错误物体就算错误。摘要和结论将按数据集分别写出最高的 Acc@0.5m：NR3D 上是带 NS-Grounding 的 Qwen3-plus（58.8），SR3D 上是带 NS-Grounding 的 GLM-5（67.0）。}

\noindent 更正后的数字见表~\ref{tab:csvg} 与表~\ref{tab:ablation}。

\subsection*{关于 between 关系}

\revcomment{论文把查询定义成二元边构成的图，但 \emph{between} 同时依赖于一个目标和两个参照物体。它是三元约束，而不是图 5 所画的两条独立二元关系。请修改表述，使图、公式和实现一致。}

\revresponse{审稿人对这一不一致的表述是准确的，我们感谢这一精确指出。实现中 \emph{between} 已经按一个源实例加两个参照实例来打分：分数是这三个自变量的函数。正文仍把查询画成二元边构成的图，图 5 仍画出两个分开的 between 标记。我们将把书面因子改成一个三元项，并重绘图 5，使一个 between 因子落在候选目标和两个锚点上。打分器本身不改。这条意见不要求新实验。表~\ref{tab:scorers} 中“去掉 between”一行是对应的实证：去掉这一三元项使 T-Acc@0.25 下降 2.4 个百分点，使 T-Acc@0.50 下降 4.4 个百分点。}

\revchange{between 分数是三个实例的函数：候选目标以及两个被指称的锚点。它不是两条二元边之和。图 5 在这三个结点上画出一个 between 因子，与实现一致。}

\subsection*{关于图 1，以及 streaming 与 online}

\revcomment{图 1 把 OnlineAnySeg 标成离线，而表 I 把它标成 Online。streaming 与 online 没有定义。}

\revresponse{修改稿图 1 不再给 OnlineAnySeg 使用离线标签。该图与表 I 一致，OnlineAnySeg 为 Online。streaming 指由 RGB-D 序列增量融合。online 指地图在序列进行中更新，而不是整段结束后再做一次离线重建。见修改稿正文图 1。}

\revchange{OnlineAnySeg 在表 I 和图 1 中都是 Online。streaming 是 RGB-D 序列的增量融合。online 指地图在序列进行中更新。}

\subsection*{关于参考文献与笔误}

\revcomment{若干方法和模型未引用，文献 [23] 格式有误，文献 [3]、[10]、[17] 已有正式发表版本。列出的笔误应改正。}

\revresponse{修改稿补充了 ZSVG3D、CSVG、审稿意见中的场景图与建图文献，以及 FCGF、体素哈希 TSDF、TSDF fusion、Qwen3-plus 和 deepseek-v3.2。GLM-5 不再引用 2022 年的 GLM 预训练论文。式 (9) 引用属性图匹配文献，式子本身仍是加权几何平均。文献 [3]、[10]、[17] 只有在能够核对为同一篇论文时才换成已发表版本。对不上 ICLR 2025、RSS 2023 或 ICCV 2025 的，不猜测配对。已指出的笔误在正文改正。}

\revchange{正文使用的文献条目都是能够核对的条目。Qwen3-plus 与 deepseek-v3.2 按解析查询图所用的模型引用。}

\bigskip
\noindent 再次感谢审稿人。这些意见促成了重建因素的分离、指标定义、CSVG 比较、绑定比较，以及更正后的定位数字。

\end{CJK*}
\end{document}
